\chapter{Hohenberg--Kohn and Kohn--Sham Theory}
\label{ch:hkks}

\section{Functional calculus in twenty minutes}

DFT is written in the language of \emph{functionals} --- machines that eat a
whole function and return a number. The total energy $E[n]$ is a functional
of the density. You need exactly two tools.

\paragraph{Functional derivative.} For a functional $F[n]$, the functional
derivative $\fdv{F}{n(\rr)}$ is defined by the first-order response to an
arbitrary small perturbation $\delta n(\rr)$:
\begin{equation}
F[n + \delta n] - F[n]
= \int \fdv{F}{n(\rr)}\,\delta n(\rr)\, \dd\rr + O(\delta n^2).
\label{eq:funcderivdef}
\end{equation}
It is the continuum analogue of the gradient: if $n$ were a vector with
components $n_i$ (which it literally becomes on your real-space grid!), then
$\delta F = \sum_i (\partial F/\partial n_i)\,\delta n_i$, and
$\fdv{F}{n(\rr)}$ is just $\partial F/\partial n_i$ divided by the grid-cell
volume. Three examples you will use constantly:
\begin{align}
F[n] &= \int \vext(\rr)\,n(\rr)\,\dd\rr
&&\Longrightarrow&
\fdv{F}{n(\rr)} &= \vext(\rr),
\\[2pt]
E_{\mathrm H}[n] &= \tfrac12 \iint
\frac{n(\rr)\,n(\rr')}{|\rr - \rr'|}\, \dd\rr\, \dd\rr'
&&\Longrightarrow&
\fdv{E_{\mathrm H}}{n(\rr)} &= \int \frac{n(\rr')}{|\rr-\rr'|}\,\dd\rr'
\equiv \vH(\rr),
\label{eq:hartreederiv}
\\[2pt]
F[n] &= \int f\big(n(\rr), \nabla n(\rr)\big)\, \dd\rr
&&\Longrightarrow&
\fdv{F}{n(\rr)} &= \pdv{f}{n} - \nabla\cdot\pdv{f}{(\nabla n)},
\label{eq:gradientfd}
\end{align}
the last being the Euler--Lagrange rule (integrate by parts once), needed for
gradient-corrected functionals in Chapter~\ref{ch:xc}.

\paragraph{Constrained minimization.} To minimize $E[n]$ subject to
$\int n = N$, introduce a Lagrange multiplier $\mu$ and demand
stationarity of $E[n] - \mu\big(\int n\,\dd\rr - N\big)$:
\begin{equation}
\fdv{E}{n(\rr)} = \mu \qquad \text{for all } \rr.
\label{eq:eulerlagrange}
\end{equation}
The multiplier $\mu$ turns out to be the chemical potential (Fermi level).

\section{The Hohenberg--Kohn theorems}

Consider $N$ electrons with the universal pieces $\hat T_e + \hat V_{ee}$
plus an external potential $\vext$. Different materials differ \emph{only}
through $\vext$ (and $N$). The 1964 theorems of Hohenberg and Kohn establish
that the ground-state density is a legitimate fundamental variable.

\begin{theorem}[HK1: the density determines the potential]
The ground-state density $n_0(\rr)$ of an interacting electron system
determines the external potential $\vext(\rr)$ uniquely, up to an additive
constant. Consequently $n_0$ determines the Hamiltonian, hence \emph{all}
ground- and excited-state properties of the system.
\end{theorem}

\begin{proof}
By contradiction, in two steps. Assume for simplicity non-degenerate ground
states. Suppose two potentials $\vext^{(1)} \neq \vext^{(2)} + \mathrm{const}$
give ground states $\Phi_1 \neq \Phi_2$ with the \emph{same} density $n_0$.
(That $\Phi_1\neq\Phi_2$ follows because they satisfy Schr\"odinger equations
with potentials differing by more than a constant.) By the variational
principle, using $\Phi_2$ as a trial state for Hamiltonian $\Ha_1$:
\begin{equation}
E_1 < \braket{\Phi_2|\Ha_1|\Phi_2}
= \braket{\Phi_2|\Ha_2|\Phi_2}
+ \braket{\Phi_2|\hat V_{\!1} - \hat V_{\!2}|\Phi_2}
= E_2 + \int \big[\vext^{(1)} - \vext^{(2)}\big]\, n_0(\rr)\, \dd\rr .
\end{equation}
Swapping labels $1\leftrightarrow 2$:
\begin{equation}
E_2 < E_1 + \int \big[\vext^{(2)} - \vext^{(1)}\big]\, n_0(\rr)\, \dd\rr .
\end{equation}
Adding the two inequalities gives $E_1 + E_2 < E_1 + E_2$, a contradiction.
Hence two essentially different potentials cannot share a ground-state
density: the map $n_0 \mapsto \vext$ is single-valued.
\end{proof}

\begin{theorem}[HK2: variational principle for the density]
There exists a universal functional $F_{\mathrm{HK}}[n]$ (the same for every
material) such that
\begin{equation}
E_v[n] \equiv F_{\mathrm{HK}}[n] + \int \vext(\rr)\, n(\rr)\, \dd\rr
\end{equation}
satisfies $E_v[n] \geq E_0$ for every admissible density with $\int n = N$,
with equality if and only if $n = n_0$ (for a non-degenerate ground state;
with degeneracy, equality holds for any ground-state density).
\end{theorem}

\begin{proof}
The cleanest construction is Levy's constrained search. Define
\begin{equation}
F_{\mathrm{HK}}[n] \equiv
\min_{\Phi \to n} \braket{\Phi | \hat T_e + \hat V_{ee} | \Phi},
\label{eq:levysearch}
\end{equation}
the minimum over all antisymmetric $N$-electron wavefunctions that yield the
prescribed density $n$. This is manifestly universal: $\vext$ does not appear.
Then for any $n$,
\begin{equation}
E_v[n]
= \min_{\Phi\to n}\braket{\Phi|\hat T_e + \hat V_{ee}|\Phi}
+ \int \vext\, n
= \min_{\Phi\to n}\braket{\Phi|\Ha_{\mathrm{el}}|\Phi}
\;\geq\; \min_{\Phi}\braket{\Phi|\Ha_{\mathrm{el}}|\Phi} = E_0 ,
\end{equation}
because searching within a fixed density is more restrictive than searching
over all wavefunctions. Minimizing over $n$ removes the restriction, so the
minimum of $E_v[n]$ is $E_0$, attained at $n_0$.
\end{proof}

\begin{remark}
These theorems are an existence statement, not a recipe: nobody knows
$F_{\mathrm{HK}}[n]$ in closed form --- knowing it exactly would be
equivalent to solving the many-body problem. The genius of Kohn and Sham
(1965) was to isolate the \emph{large, known} parts of $F_{\mathrm{HK}}$ and
approximate only the small remainder.
\end{remark}

\section{The Kohn--Sham construction}
\label{sec:ks}

The dominant contribution to the kinetic energy is well approximated by that
of \emph{non-interacting} electrons with the same density. So: introduce an
auxiliary fictitious system of $N$ non-interacting electrons moving in some
effective potential $\veff(\rr)$, engineered so that its ground-state
density equals the true interacting density. Non-interacting electrons are
easy: they occupy single-particle orbitals $\psi_i(\rr)$ (for now spinless;
spinors come in Chapter~\ref{ch:soc}) obeying a one-body Schr\"odinger
equation, and
\begin{equation}
n(\rr) = \sum_{i} f_i\, |\psi_i(\rr)|^2 ,
\qquad
T_s[n] = \sum_i f_i \int \psi_i^*(\rr)
\Big(-\tfrac{1}{2}\nabla^2\Big) \psi_i(\rr)\,\dd\rr,
\label{eq:ksdensity}
\end{equation}
with occupations $f_i \in [0,1]$ and $\sum_i f_i = N$. (Here $i$ counts
\emph{spin-orbitals}: one state, one electron at most. If spin is a pure
spectator one can instead let each spatial orbital hold two electrons; we
will not, because with spin--orbit coupling every state is its own spinor.)
The many-electron wavefunction of non-interacting fermions is a single
\emph{Slater determinant} of the occupied orbitals,
$\Phi = (N!)^{-1/2}\det[\psi_i(\rr_j)]$ --- the antisymmetrized product
``one electron in each orbital''; Eq.~\eqref{eq:ksdensity} is
what Eq.~\eqref{eq:densitydef} gives for a determinant.
$T_s$ is the \emph{non-interacting} kinetic energy --- not the true one, but
its largest part, and crucially one we can evaluate exactly through the
orbitals.

Now \emph{define} the exchange--correlation functional as everything left
over:
\begin{equation}
\boxed{\;
F_{\mathrm{HK}}[n] = T_s[n] + E_{\mathrm H}[n] + \Exc[n]
\;}
\qquad\Longleftrightarrow\qquad
\Exc[n] \equiv \big(T - T_s\big) + \big(V_{ee} - E_{\mathrm H}\big).
\label{eq:excdef}
\end{equation}
$\Exc$ contains (i) the difference between true and non-interacting kinetic
energies, (ii) the exchange energy --- the energy lowering from Pauli
antisymmetry keeping same-spin electrons apart, (iii) the correlation energy
--- the further lowering from electrons dynamically dodging one another
beyond mean field, and (iv) the correction for the self-interaction hiding
inside $E_{\mathrm H}$ (in $E_{\mathrm H}$ each electron feels its own
charge). $\Exc$ is a small fraction of the total energy ($\sim$1--10\%), yet
it contains all the chemistry; approximating it is Chapter~\ref{ch:xc}.

\section{Deriving the Kohn--Sham equations}

Apply the Euler--Lagrange condition \eqref{eq:eulerlagrange} to
$E_v[n] = T_s[n] + E_{\mathrm H}[n] + \Exc[n] + \int \vext n$:
\begin{equation}
\fdv{T_s}{n(\rr)} + \vH(\rr) + \vxc(\rr) + \vext(\rr) = \mu ,
\qquad
\vxc(\rr) \equiv \fdv{\Exc}{n(\rr)} .
\label{eq:eltrue}
\end{equation}
Now write the same stationarity condition for the fictitious non-interacting
system in potential $\veff$:
\begin{equation}
\fdv{T_s}{n(\rr)} + \veff(\rr) = \mu' .
\label{eq:elks}
\end{equation}
Equations \eqref{eq:eltrue} and \eqref{eq:elks} produce the same density if
and only if
\begin{equation}
\boxed{\;
\veff(\rr) = \vext(\rr) + \vH[n](\rr) + \vxc[n](\rr)
\;}
\label{eq:veff}
\end{equation}
(up to a constant). But the ground state of non-interacting electrons in
$\veff$ is found by solving the one-body eigenproblem. Minimizing
$E$ with respect to the orbitals themselves, subject to orthonormality
$\braket{\psi_i|\psi_j} = \delta_{ij}$ (Lagrange multipliers
$\varepsilon_{ij}$, diagonalized to $\varepsilon_i$), gives the same result
directly. A complex function has two independent real degrees of freedom,
its real and imaginary parts; it is equivalent, and shorter, to treat
$\psi_i$ and $\psi_i^*$ as the two independent variables and set the
derivative with respect to $\psi_i^*$ to zero (the derivative with respect
to $\psi_i$ gives the complex-conjugate equation). Varying $\psi_i^*$,
\begin{equation}
\fdv{}{\psi_i^*(\rr)}\Big[E - \textstyle\sum_{jk}
\varepsilon_{jk}\big(\braket{\psi_j|\psi_k} - \delta_{jk}\big)\Big] = 0
\;\Longrightarrow\;
\Big[-\tfrac12 \nabla^2 + \veff(\rr)\Big] \psi_i(\rr)
= \varepsilon_i\, \psi_i(\rr).
\label{eq:kseq}
\end{equation}
(Strictly, varying $\psi_i^*$ gives
$f_i\,\hat h\,\psi_i = \sum_k \varepsilon_{ik}\psi_k$ with
$\hat h = -\tfrac12\nabla^2+\veff$, because $\psi_i$ enters $E$ with weight
$f_i$. Project onto $\psi_j$: $f_i\braket{\psi_j|\hat h|\psi_i} =
\varepsilon_{ij}$. The multiplier matrix can be taken Hermitian, and so is
$\hat h$; combining the $(i,j)$ and $(j,i)$ equations gives
$(f_i - f_j)\braket{\psi_j|\hat h|\psi_i} = 0$. So $\hat h$ has no matrix
elements between orbitals of different occupation, and within a set of
equally occupied orbitals a unitary rotation --- which changes neither $n$
nor $E$ --- diagonalizes it. Hence the form above.)
Equations \eqref{eq:ksdensity}, \eqref{eq:veff} and \eqref{eq:kseq} are the
\textbf{Kohn--Sham equations}. Note the nonlinearity: $\veff$ depends on $n$,
which depends on the $\psi_i$, which depend on $\veff$. The equations must be
solved \emph{self-consistently} --- the fixed-point iteration at the heart of
your program (Chapter~\ref{ch:algorithms}).

\section{The total energy}

Given a self-consistent solution, the ground-state energy is
\begin{equation}
E = T_s + \int \vext\, n\, \dd\rr + E_{\mathrm H}[n] + \Exc[n] + E_{NN}.
\label{eq:etot1}
\end{equation}
A form better suited to computation eliminates $T_s$ using the eigenvalues.
Multiply \eqref{eq:kseq} by $f_i\psi_i^*$, integrate, and sum:
$\sum_i f_i \varepsilon_i = T_s + \int \veff\, n$. Hence
\begin{equation}
\boxed{\;
E = \sum_i f_i\,\varepsilon_i
\;-\; E_{\mathrm H}[n]
\;-\; \int \vxc(\rr)\, n(\rr)\, \dd\rr
\;+\; \Exc[n]
\;+\; E_{NN}.
\;}
\label{eq:etot2}
\end{equation}
The corrections after the band-structure sum $\sum f_i\varepsilon_i$ are
called \emph{double-counting} terms: the eigenvalue sum counts the Hartree
energy twice (each electron feels the full potential of all electrons) and
counts $\int \vxc n$ instead of $\Exc$. For a metal with smeared occupations
an entropy term is added (Section~\ref{sec:smearing}). Both forms are
worth coding: Eq.~\eqref{eq:etot1} (``direct'') and Eq.~\eqref{eq:etot2}
(``eigenvalue sum'') must agree to rounding error at self-consistency, and
a disagreement pinpoints a bookkeeping bug.

\begin{keybox}{Interpretation warnings}
The Kohn--Sham orbitals and eigenvalues belong to the \emph{fictitious}
system. Their only rigorous roles are to build $n(\rr)$ and $T_s$.
Nevertheless, the KS band structure of metals like gold is a good
approximation to the measured quasiparticle bands (photoemission; for gold
the $d$ bands come out several tenths of an eV too shallow --- $0.3$--$0.9$\,eV
with the functionals and pseudopotential of this document,
Chapters~\ref{ch:soc} and~\ref{ch:colour}), and
(for the exact functional) the highest occupied eigenvalue is exact: it
equals minus the ionization energy of a finite system, and the Fermi level
of a metal. Band
\emph{gaps} of semiconductors are systematically underestimated --- a famous
limitation that does not affect our metallic target.
\end{keybox}

\section{Spin-density functional theory}

When a magnetic field couples to spin, or when SOC entangles spin and orbit,
the theory is generalized (von Barth--Hedin 1972): the energy becomes a
functional of the full spin-density matrix
$n_{\alpha\beta}(\rr)$ of Eq.~\eqref{eq:spindensmat},
$E = E[n_{\alpha\beta}]$, and the XC potential becomes a $2\times2$ matrix
field,
\begin{equation}
v^{\mathrm{xc}}_{\alpha\beta}(\rr)
= \fdv{\Exc}{n_{\beta\alpha}(\rr)}
= \vxc(\rr)\,\delta_{\alpha\beta}
+ \mathbf{B}_{\mathrm{xc}}(\rr)\cdot\bm\sigma_{\alpha\beta},
\label{eq:vxcmatrix}
\end{equation}
with an \emph{exchange--correlation magnetic field}
$\mathbf{B}_{\mathrm{xc}} = \delta \Exc / \delta \mathbf{m}$. The Kohn--Sham
orbitals become two-component spinors
$\psi_i = (\psi_i^{\uparrow}, \psi_i^{\downarrow})^{\mathsf T}$ and the KS
equation a $2\times 2$ matrix equation --- exactly the structure SOC requires,
developed fully in Chapter~\ref{ch:soc}. For non-magnetic gold,
$\mathbf{m} = 0$ and $\mathbf{B}_{\mathrm{xc}} = 0$ at self-consistency, but
your data structures should support the general case.
